\documentclass[11pt,a4paper]{article}
\usepackage[margin=2.4cm]{geometry}
\usepackage[T1]{fontenc}
\usepackage[utf8]{inputenc}
\usepackage{lmodern,microtype,graphicx,booktabs,longtable,tabularx,enumitem,xcolor,url}
\usepackage[round,authoryear]{natbib}
\usepackage[hidelinks]{hyperref}
\graphicspath{{images/}{figures/}}
\title{C-LARA-2: Initial Project Report\\AI Autonomy Under Human Supervision in a Multimodal CALL Platform}
\author{%
\parbox{0.95\textwidth}{\centering
\normalsize Branislav B{\'e}di\textsuperscript{1}, ChatGPT C-LARA-Instance\textsuperscript{2}, Belinda Chiera\textsuperscript{3}, and Cathy Chua\textsuperscript{4}\\
\normalsize St{\'e}phanie Geneix-Rabault\textsuperscript{5}, Christ{\`e}le Maizonniaux\textsuperscript{6}, Manny Rayner\textsuperscript{4}, and Sophie Rendina\textsuperscript{7}\\
\normalsize Sarah Wright\textsuperscript{4} and Rina Zviel-Girshin\textsuperscript{8}\\[0.65em]
\footnotesize \textsuperscript{1}The {\'A}rni Magn{\'u}sson Institute for Icelandic Studies, Iceland\\
\footnotesize \textsuperscript{2}OpenAI large language model instance participating in the C-LARA project\\
\footnotesize \textsuperscript{3}Adelaide University, Australia; \textsuperscript{4}C-LARA Project, Australia; \textsuperscript{5}University of New Caledonia, New Caledonia\\
\footnotesize \textsuperscript{6}Flinders University, Australia; \textsuperscript{7}University of New England, New South Wales, Australia\\
\footnotesize \textsuperscript{8}Ruppin Academic Center, Israel
}}
\date{July 2026}
\begin{document}
\maketitle
\phantomsection
\addcontentsline{toc}{section}{Abstract}
\begin{abstract}
C-LARA-2 is a new, AI-mediated implementation of the C-LARA platform for creating multimodal language-learning resources. This initial report records the first nine months of development and examines both the emerging platform and the development method: a small international team has used a high-capability AI coding assistant to write code, tests, documentation, roadmaps and much of the report itself, under strategic human supervision. We describe the modular linguistic-annotation pipeline, image-generation and community-review workflows, specialised picture dictionaries, repository-based project management, AWS deployment, a project-understanding Assistant and an exploratory use of Voice Mode for spoken practice. The report draws in part on two related EuroCALL papers while providing a broader, repository-grounded account of the project's current state, limitations and plans. The case studies suggest that AI assistance can greatly reduce the cost of implementation, redesign, documentation and operational work, while leaving human collaborators responsible for pedagogical purpose, community knowledge, governance and review. We conclude by considering how AI autonomy, including strategic and ethical advice, might be progressively delegated while preserving accountable human and community control.
\end{abstract}
\phantomsection
\addcontentsline{toc}{section}{Table of contents}
\tableofcontents
\clearpage
% During drafting, use e.g. \includeonly{sections/04_annotation_pipeline}
% \include starts each major section on a new page. Change to \input later if a compact layout is preferred.
\section{Introduction}

C-LARA-2 is an AI-based online platform for creating and using multimodal
language-learning materials. It is a complete reimplementation of C-LARA, a
platform developed to bring together texts, linguistic annotation,
translations, audio, images and interactive learning activities. C-LARA-2
retains those pedagogical purposes while rebuilding the platform around a
cleaner technical architecture and a distinctive development method: the
project delegates extensive operational work to an AI coding agent while
retaining human responsibility for strategy, values, evaluation and
acceptance.

The central claim of this report is consequently not that C-LARA-2 is a
finished system, nor that every component is equally mature. It is that the
project is an ongoing experiment in \emph{AI autonomy under strategic human
supervision}. Human collaborators decide what the project is for, identify
learner and community needs, supply linguistic and pedagogical expertise,
set ethical constraints, assess proposed solutions and decide whether they
are acceptable. Codex takes responsibility for much of the technical work:
interpreting requirements, inspecting the repository, designing and changing
code, writing tests, maintaining documentation, organising experiments and
preparing technical and publication drafts. The division is deliberately
asymmetric. The AI has substantial freedom over operational means, but does
not acquire authority to set the project's purposes or standards.

This model matters because conventional academic software projects often
face a mismatch between the breadth of their ambitions and the development
capacity available to a small, distributed consortium. A coding agent can
make implementation and revision much faster, but speed is not itself the
point. Its value is that it can make specialised, user-led changes practical:
for example, a picture-dictionary interface can be redesigned after users
reject the first interaction model without discarding the underlying data
representation; documentation, roadmaps and issue records can be kept close
to the implementation they describe; and an evolving platform can be
re-examined and redirected at relatively low operational cost. These gains
also create a verification burden. Producing code, documentation and prose
quickly does not make them correct. The project therefore treats version
control, tests, inspectable documentation, issue records, human review and
reversible changes as necessary counterparts to operational autonomy.

The accelerated pace of AI-assisted development also creates a problem for
academic communication. When an AI agent can support implementation,
documentation, analysis and drafting in rapid cycles, a project can produce
substantially more work in the time and with the resources that would
previously have supported a much smaller programme. The difficulty is not
only that conventional publication takes time; it is that a detailed
description of a rapidly changing system can become historically accurate
but practically outdated before it appears.

This pressure reflects a broader tension between established academic
publication practices and AI-mediated research and development. \citet{Manabu2026} argues that generative AI may undermine a scholarly economy
organised around the production of countable publication artefacts, while
\citet{KobakEtAl2025} provide evidence that LLM-assisted scholarly writing
has already become widespread. We do not attempt to decide whether these
developments will destroy the familiar ``publish-or-perish'' model. Our
more conservative claim is that they are already changing the relation
between the pace of research work and the pace at which that work can be
publicly documented, evaluated and credited.

C-LARA-2 therefore presents its initial work in two complementary forms.
A short paper submitted to EuroCALL 2026 in July 2026 will be published
in December 2026 and can function as a conventional peer-reviewed
publication. Because the platform will continue to change during that
interval, the EuroCALL paper concentrates on the project's philosophical
foundations and uses a small number of illustrative examples rather than
attempting a comprehensive functional account. The present report takes
the opposite approach. It records the July 2026 platform in substantially
greater detail, including implemented areas, selected experiments and
current limitations, and will be posted immediately on ResearchGate.
This gives the report the status of a current, citable project snapshot,
but not the academic prestige or independent validation attached to peer
review.

The two forms should not be confused. The conference paper offers a
selective, externally reviewed account of the development model; this
report offers a fuller and more rapidly available record of the system at
a defined point in its history. The need to combine them is itself
evidence that the questions raised by Manabu have practical force. It
does not settle those questions, but it illustrates how a small academic
project can respond when AI-assisted work evolves more quickly than the
ordinary publication cycle.

The report records selected episodes that make the development model
concrete: the annotation and image-generation pipelines, picture
dictionaries and the learning activities derived from them, project
management and deployment work, and the repository-grounded Assistant
that can answer questions about the platform while distinguishing
implemented behaviour from planned work. These examples are evidence
about a method of organising work, not proof that the platform or its AI
processes are complete, autonomous in every sense, or free of error.

A further strand concerns spoken conversational practice. Sarah Wright has
conducted regular sessions using ChatGPT Pro Tier Voice Mode as a conversation
partner. These sessions provide a practically important contrast and
complement to C-LARA-2's text-centred multimodal resources: a learner can
work with a prepared text and then sustain spoken interaction around it with
a responsive AI partner. The resulting experience is promising, but the two
environments are not yet naturally integrated. The report will document what
these sessions suggest about the pedagogical opportunity, the present
workarounds, and the interface and evidence questions that must be addressed
before any stronger claim is made about an integrated conversational
functionality.

The remainder of the report is organised as follows. Section~2 describes the
history of C-LARA and C-LARA-2 and explains the human/AI development method.
Section~3 gives a platform overview. Sections~4--6 cover the annotation
pipeline, image generation and picture dictionaries; Section~7 addresses
project management; and Section~8 records deployment work. Section~9
examines the project-understanding Assistant, while Section~10 develops the
Voice Mode and conversational-practice thread. Section~11 identifies future
directions, including stronger evaluation and integration work. Finally,
Section~12 draws together the report's conclusions about the benefits,
limits and governance requirements of operational AI autonomy under strategic
human supervision.

\section{History and methodology}

\subsection{C-LARA: pedagogical aims and workflow}

C-LARA, the ChatGPT-Based Learning And Reading Assistant, was developed
shortly after the release of GPT-4 in March 2023. Its purpose was to make it
practical to create multimodal language-learning materials for a wide range
of languages and learners. A typical project starts with a text supplied by a
user or generated from a short specification. The text is then enriched with
interconnected linguistic and pedagogical resources: page and segment
structure, lexical segmentation, translations, multi-word expressions,
lemmas, parts of speech, glosses, audio and illustrations. The resulting
material can be compiled as an interactive web resource in which learners
move among textual, linguistic, auditory and visual representations.

The workflow was designed as an end-to-end authoring route rather than a
collection of isolated tools. A user can create or import text; invoke
successive annotation operations, either as a sequence or with review between
stages; create page-aligned illustrations; and deploy the compiled resource
as interactive HTML. Annotation operations include segmentation, segment
translation, MWE identification, lemma/POS tagging, glossing and audio
attachment. Image creation can use a style description, reusable visual
elements and page-level images. The details have evolved, but the pedagogical
principle remains stable: AI services should make it easier for teachers and
other domain experts to create, inspect and revise materials rather than
replace their judgement \citep{ChatGPTetal2023,RendinaEtAl2025}.

\subsection{Experience across languages and communities}

C-LARA has been used with both widely taught and low-resource languages.
Earlier work evaluated GPT-4 processing for English, Farsi, Faroese, Mandarin
and Russian \citep{ChatGPTetal2023}. Work involving the Kanak languages
Drehu and Iaai combined native-speaker audio, AI-generated illustrations,
linguistic annotation and culturally grounded content
\citep{BaeEtAl2025,ChieraEtAl2025,WelbyEtAl2024}. These examples do not show
that AI processing is equally reliable for every language. In particular,
where the project language is not well supported by an available model,
human-provided translations in a language such as English or French may
mediate some operations, notably image generation. Community knowledge,
cultural permission, data sovereignty, linguistic judgement and audio remain
human responsibilities.

The history is important because it established both the opportunity and the
limit of the platform. A small consortium can create useful multimodal
materials for specialised communities when AI reduces routine operational
work. But the value of the result depends on human collaborators who know the
learners, language and community context well enough to set requirements and
assess outcomes.

\subsection{Why a rewrite was necessary}

The original implementation grew incrementally over several years. GPT-4 and
later ChatGPT models wrote a substantial proportion of the code under close
human supervision; the project's informal estimate is approximately 75\%.
That AI-assisted model made a broad platform feasible for a small academic
consortium. It also left the familiar consequences of repeated local
extension: accumulated architectural compromises, uneven documentation and
components that became difficult to understand or maintain.

C-LARA-2 began in November 2025 as a rational rewrite. The aim was not to
abandon C-LARA's pedagogical model, but to preserve its successful concepts
in a cleaner and more systematic technical architecture. For several months,
the two systems were developed in parallel while the project compared them.
In March 2026, development on C-LARA was halted after C-LARA-2 had become the
clearly preferable platform. The rewrite was thus also an opportunity to
reconsider how the technical representation of the project itself could be
maintained.

\subsection{From AI-assisted development to operational AI autonomy}

The availability of Codex through the ChatGPT Pro Tier changed the nature of
that opportunity. Instead of using an AI principally as a closely supervised
source of implementation suggestions, C-LARA-2 assigned it broad operational
responsibility for the evolving repository. At the time of this report, Codex
has produced all committed program code, tests and repository documentation.
It also maintains roadmaps, issue records, technical explanations, operational
notes, experiment organisation and publication drafts. Human collaborators
provide requirements, examples, linguistic and pedagogical expertise,
strategic priorities, ethical constraints, review and acceptance decisions;
they do not directly author the implementation.

This distinction is central. C-LARA-2 is not simply a conventionally written
system accelerated by code completion. Code, tests, documentation, roadmaps
and issue records are maintained together as a shared project memory that can
be inspected by both the human collaborators and the AI agent. The project is
therefore testing whether an AI can assume continuing responsibility for the
operational representation of a complex CALL platform while humans retain
strategic authority over its purposes, values and standards. This delegation
does not eliminate the need for verification; it makes version control,
testing, inspectable documentation, human review and reversible change more
important.

The next section describes the July 2026 C-LARA-2 platform that resulted from
this history and organisational model.

\section{Platform overview and report scope}

This report describes the July 2026 version of C-LARA-2 as a current project
snapshot. It draws on the implemented platform, repository source code and
tests, roadmap and issue records, operational documentation, selected
experiments and human collaborator experience. The aim is not to provide a
controlled evaluation of every capability, but to give an inspectable account
of the platform's principal functional areas, the development model through
which they are being extended, and the limitations that remain.

At the highest level, C-LARA-2 supports the construction and use of
multimodal language-learning projects. Users can create or import texts;
annotate them with linguistic, translation and audio information; generate
and select images; compile the results as interactive web resources; and
manage projects and their versions. The platform retains the C-LARA idea that
these activities are connected stages in the production of learner-facing
material, while providing more systematic support for specialised workflows,
revision and operational visibility.

The first group of functions concerns text creation and linguistic
processing. A user can begin with source text or an AI-generated text, choose
relevant project and glossing languages, and run an annotation pipeline that
includes segmentation, translation, MWE analysis, lemma/POS tagging and
glossing. The pipeline is not treated as an infallible black box: users can
review and edit outputs, and the repository retains prompts, tests,
experiments and issue records through which processing behaviour can be
examined and improved.

A second group concerns multimodal presentation. Image-generation workflows
can produce page illustrations and support consistent visual style; picture
dictionaries adapt the underlying project representation to a task in which a
lexical entry, its translation, its image-generation prompt and its selected
image should be managed together. The same resources can support picture
glossing and picture-based exercises. These functions are especially relevant
where accessible visual materials can help learners or community contributors,
but their use remains subject to human pedagogical review and, where
appropriate, cultural constraints.

A third group concerns the infrastructure that makes an evolving platform
usable and inspectable. Project-management features, deployment work,
roadmaps, issue records and operational documentation make it possible to
track work beyond an individual interaction with an AI agent. The
project-understanding Assistant builds on this repository memory: it can
inspect current project artefacts and answer natural-language questions about
implemented behaviour, planned work and design decisions, with links back to
the supporting repository evidence. This is an early, limited form of project
self-understanding rather than a claim of perfect knowledge or autonomous
judgement.

Finally, C-LARA-2 is beginning to explore how prepared multimodal texts might
relate to spoken conversational practice. Sarah Wright's regular sessions with
ChatGPT Pro Tier Voice Mode illustrate a promising conversational complement
to the platform's text-centred resources. The two environments are not yet
naturally integrated, and the report distinguishes present workarounds from
implemented C-LARA-2 functionality.

Sections~4--10 examine these areas in more detail: the annotation pipeline,
image generation, picture dictionaries, project management, deployment, the
project-understanding Assistant and Voice Mode. The organisation is
deliberately selective. It foregrounds functions and episodes that clarify
the C-LARA-2 development model, rather than attempting a complete feature
manual. Sections~11 and~12 then identify the most important unfinished work
and draw together the report's conclusions about operational AI autonomy under
strategic human supervision.

\section{Annotation pipeline}
\todo{Draft this section.}

\section{Image generation}

\subsection{Purpose and workflow}

C-LARA-2 retains the image-generation workflow developed in the final versions
of C-LARA. Its purpose is not merely to attach an attractive image to each
page. A multi-page pedagogical text needs illustrations that fit the content
of individual pages while also forming a coherent visual world: recurring
people, objects and locations should remain recognisable, and linework,
colour, tone and degree of realism should not vary arbitrarily. The
availability of GPT-Image-1 to the project in 2025 made this workflow much
more practical than it had been with earlier image-generation models
\citep{RendinaEtAl2025}.

The workflow separates three reviewable stages. First, a human author supplies
a short style brief. The system expands it into a fuller description and
creates a sample image; the author can inspect, edit or regenerate both before
continuing. Second, the system proposes recurring visual \emph{elements}---for
example, people, objects and locations. The author can remove unsuitable
proposals or add missing ones, and the platform creates descriptions and
reference images for the selected elements. Finally, it creates one or more
candidate page images using the page text, approved style, relevant element
references and any page-level guidance. Independent element and page requests
can be run in parallel, but selection and acceptance remain human decisions.

\subsection{A two-page French example}

We illustrate the workflow with the small French dialogue already used in the
annotation discussion:

\begin{quote}
\textbf{Page 1:} «~Je t'aime~», dit-il.\\
\textbf{Page 2:} «~Vous vous moquez de moi~», dit-elle.
\end{quote}

The example is deliberately minimal. Its purpose is not to claim literary or
pedagogical depth, but to make the relationship among style, elements and
page-level generation visible. The two lines create a small unresolved story:
a declaration is followed by a response suggesting that the woman believes the
man is making fun of her. This gives the generator a local emotional task on
each page while requiring continuity of character and visual atmosphere.

For the style stage, the author supplied the brief ``Late nineteenth-century
French lithograph-inspired illustration, with restrained hand-coloured tones,
fine ink linework, soft paper texture, and a slightly melancholic literary
atmosphere.'' The brief also requested suitability for adult or older teenage
language learners and prohibited visible written text, speech balloons, labels
and signage. The platform expanded this in French into a more detailed prompt
specifying delicate ink lines, muted hand-coloured tones, paper texture, soft
light, subdued shadows and an intimate, contemplative composition. This is a
useful example of the division of labour: the human supplied the historical
and pedagogical intention, while the AI converted it into a fuller operational
specification that remained available for human review.

\subsection{Recurring elements and page images}

Element discovery proposed \emph{l'homme} and \emph{la femme}. Their expanded
descriptions specify nineteenth-century dress, fine ink detail, muted colour,
soft lighting and emotionally expressive but restrained faces. The resulting
reference images establish a recognisable visual identity for each person.
They are then provided as relevant element references for both page prompts,
so that the system has an explicit continuity mechanism rather than relying
only on repeated textual descriptions.

The page prompts are intentionally close to the element and style
specifications. Page 1 foregrounds the man and the emotionally tentative
moment of the declaration; Page 2 foregrounds the woman and her more guarded
response. Both prompts retain references to the woman and the man, even where
only one is foregrounded. Three variants were requested for each page, after
which a human author selected the preferred variant. This is an important
practical point: generation produces candidates, not an automatically accepted
final illustration.

% Repository figure files remain text-only placeholders. In Overleaf, replace
% them with a montage of the selected element-reference and page-image JPEGs.
% \begin{figure}[htbp]
% \centering
% \includegraphics[width=\linewidth]{images/lithograph_elements.jpg}
% \caption{Reference images for \emph{l'homme} and \emph{la femme}, generated
% from the approved late-nineteenth-century French lithograph style.}
% \label{fig:lithograph-elements}
% \end{figure}
%
% \begin{figure}[htbp]
% \centering
% \includegraphics[width=\linewidth]{images/lithograph_pages.jpg}
% \caption{Selected page-image variants for the two-page French dialogue. The
% common style and element references aim to preserve visual continuity while
% allowing the emotional focus to change between pages.}
% \label{fig:lithograph-pages}
% \end{figure}

\subsection{Human review and limitations}

The example illustrates why image generation is organised as a staged
workflow. A single page prompt can produce a plausible isolated picture, but
it gives little control over cross-page consistency. Separating style,
recurring elements and page images gives a human author several points at
which to inspect and redirect the process. It also makes failures easier to
locate: a problem may concern the style description, an element identity, a
page prompt or the selection among generated variants.

The workflow does not make visual or cultural judgement automatic. Generated
images can still misrepresent a relationship, omit a relevant detail, produce
inconsistent character appearances or be inappropriate for a learner or
community context. Visible text in images may also be actively undesirable,
particularly for picture-based exercises. For this reason, C-LARA-2 treats
image generation as AI-supported production under human review, with editing,
regeneration and selection available throughout. These concerns become still
more important for the specialised picture-dictionary workflow discussed in
the next section.

\section{Picture dictionaries and multimodal practice}

\subsection{Why picture dictionaries matter}

Picture dictionaries are a useful point of contact between C-LARA-2's
multimodal authoring infrastructure and the needs of minority and Indigenous
language communities. Picture-based resources can support learners with
limited literacy, communities whose orthographies are still being consolidated,
and contributors who want to create accessible resources for younger learners.
They also make a demanding requirement especially clear: community members
must be able to author and review individual lexical entries without having to
master the technical workflow that underlies a general-purpose annotated text.

The original technical motivation for adding picture dictionaries was picture
glossing. In this learner-facing functionality, a word or expression in a
C-LARA-2 text is associated with a dictionary entry; clicking it can display
the corresponding picture alongside other support such as audio playback and
concordance information. The required infrastructure is a reviewed, reusable
association between lexical items and images. Picture glossing has now been
implemented, but has not yet been tested in learner use.

In practice, the first substantial demand for picture dictionaries has come
from collaborators working in Indigenous community contexts, where the
immediate use cases are not identical to the original picture-glossing
scenario. Contributors need accessible ways to create, inspect and control
picture-based lexical resources and to reuse them in exercises. These needs
made the dictionary-entry workflow, including visible cultural guidance and
human review of images, the more urgent design problem.

A picture dictionary is represented internally as an ordinary C-LARA-2 project
whose text has one lexical item or lexicalised expression per page. This was a
useful reuse decision: it made existing project, annotation, image and
compilation infrastructure available without introducing a separate data
model. It did not, however, imply that the standard project interface was the
right interface for dictionary construction.

\subsection{Why the general-purpose workflow was not enough}

The annotation interface described in Section~4 is designed for an extended
text containing pages, segments and token-level linguistic analysis. The image
workflow in Section~5 is designed for page-level scenes whose style and
recurring visual elements must be coordinated. Both workflows are appropriate
for their intended tasks. A picture dictionary is more constrained: the
relevant unit is a lexical entry, which needs a word, perhaps a lemma and part
of speech, a gloss or translation, optional human guidance, an image prompt
and a selected image considered together.

The initial picture-dictionary design exposed the underlying project as a
sequence of existing annotation and page-image operations. Contributors had to
move between separate views to manage information that conceptually belonged
to one entry. This created complexity without a corresponding benefit,
particularly where contributors were community members rather than technical
specialists, or where the project language could not itself be processed
reliably by the available AI model. In such cases, a human-provided English or
French translation may be needed to mediate prompt creation, while cultural
notes or restrictions must remain visible to the person reviewing the image.

\subsection{A dictionary-entry-centred interface}

Following discussion with intended users, C-LARA-2 retained the project-based
representation but replaced the user-experience layer with a specialised
picture-dictionary editor. Each lexical entry is displayed as an integrated
block containing the word, lemma, part of speech, gloss or translation,
optional suggestion, image-generation prompt and selected image. Contributors
can inspect and edit each component directly; they can also select entries in
bulk to create missing translations, prompts, images or other information.
The interface thus exposes the relevant parts of the annotation and image
workflows without asking a contributor to navigate the complete pipelines for
every word.

\begin{figure}[htbp]
\centering
\includegraphics[width=\linewidth]{images/picture_dictionary_editor.jpg}
\caption{Part of the specialised C-LARA-2 picture-dictionary editor. Each
lexical entry brings together word, lemma/POS information, gloss or
translation, optional human guidance, image-generation prompt and selected
image. Controls support review and batch generation for selected entries. A
French example is shown rather than community-language data because work
involving Indigenous languages may be subject to cultural-sensitivity and
data-sovereignty constraints.}
\label{fig:picture-dictionary-editor}
\end{figure}

The redesign was not merely cosmetic. The earlier interface treated dictionary
construction as a sequence of operations inherited from the general project
workflow; the replacement treats the dictionary entry as the organising unit.
Codex produced a usable implementation within a few hours once the revised
conceptual model had been agreed, and the consortium members involved in
testing responded positively. This is a useful example of AI-centred
development under human supervision: the AI did not independently identify the
community need or diagnose the usability problem; human collaborators supplied
those judgements, while Codex carried out the rapid operational redesign.

The episode also illustrates a practical consequence of AI-mediated
development. Conventional software work is subject to a sunk-cost pressure:
prior effort can make it difficult to abandon an unsuitable design
\citep{ArkesBlumer1985}. Here, the project was able to preserve the useful
underlying representation while discarding and rebuilding the user-facing
layer. The image-generation process itself still requires further work and
careful human review; rapid implementation is not evidence that generated
images are automatically pedagogically or culturally appropriate.

\subsection{Current exercises and future picture glossing}

At present, the picture-dictionary infrastructure is being used most
concretely to support picture-based exercises in Indigenous community
contexts. Picture glossing is also implemented, but should be treated as an
untested next application rather than as a demonstrated educational result.
The same community-authored lexical entries can support more than a dictionary
display. In picture glossing, clicking a word or expression in a C-LARA-2 text
can show its picture, play audio and provide concordance examples. The
platform can also create picture-to-word and word-to-picture multiple-choice
exercises, word scrambles and simple crosswords using picture clues. This
reuse matters: effort invested in reviewing a lexical entry can support
several forms of multimodal practice rather than one isolated artefact.

\begin{figure}[htbp]
\centering
\includegraphics[height=0.72\textheight]{images/word_scramble.jpg}
\caption{A picture-clue word-scramble exercise generated from a French picture
dictionary. The learner selects adjacent letters in the grid to spell the
lexical item represented by the image; two target words have already been
found.}
\label{fig:word-scramble}
\end{figure}

The word-scramble example is intentionally modest, but it makes the point
concrete. The learner receives a visual clue, performs an active lexical task,
receives feedback and moves to the next item. The relevant innovation is not
that a word scramble is novel in itself; it is that a reviewed dictionary
entry can be reused rapidly as one component of an integrated, multimodal
learning resource.

The project intends to return to the original picture-glossing idea in a
context involving multimodal texts for children or beginner learners. For
picture glossing to be useful throughout a text, a dictionary is likely to
need coverage of at least several hundred common words and expressions. Most
of the operational work can in principle be automated: candidate vocabulary
can be identified, translations and prompts proposed, candidate images
generated, and dictionary entries linked to textual tokens. Human review
remains necessary, but it can focus on accepting, correcting or rejecting
proposed entries rather than manually constructing the whole resource. This
is a promising area for pedagogical discussion with picture-based-learning
specialists, including Christ\`ele Maizonniaux, whose research interests are
closely relevant to the design choices involved.

\section{Project memory, roadmaps and issue tracking}

\subsection{Why a substantial project needs external memory}

C-LARA-2 is now a substantial evolving repository: the July 2026 snapshot
contains approximately 82,740 tracked lines accumulated over about nine
months. No participant, human or AI, can reliably retain all relevant
architectural decisions, experiments, operational procedures, unresolved
problems, user requirements and publication claims in conversational working
memory. The project therefore treats the repository as more than a container
for code. It is a version-controlled external memory that records current
behaviour, past decisions, intended direction and the procedures through which
people can inspect or change the system.

This external memory must be searchable, structured, linked to the code it
describes and maintained as the project changes. It is useful to human
collaborators, who need to understand and redirect work without becoming
full-time programmers or system administrators. It is equally useful to
Codex, which can inspect the same records when implementing a change,
diagnosing a failure, planning an experiment or preparing a report. The aim is
not to replace human judgement with accumulated text, but to make project
history and current constraints available for human review and AI-supported
operational work.

\subsection{Roadmaps and how-to documentation}

Roadmap files have been part of C-LARA-2 from the start. They record why a
feature area matters, describe current architecture and design choices,
preserve the history of false starts or revisions, identify priorities and
dependencies, and link to relevant issues, tests, commands and operational
procedures. A roadmap is consequently not a static feature wish list. It is a
maintained account of how a part of the project reached its current state and
what work should happen next.

How-to files translate that memory into procedures that human collaborators
can follow. They cover development setup, backups, deployment, server
administration, migration and diagnostics. These documents are especially
important in an AI-centred project because a human participant may need to
perform an action in a production environment while relying on Codex to
inspect repository context, explain the reason for the action and formulate
safe, concrete steps. The documents preserve that operational knowledge for
the next occasion rather than leaving it only in a transient conversation.

\subsection{Issues as AI-maintained project records}

Focused issue records were added in early May 2026. They can be entered and
updated by human collaborators through a tab in the online platform, but they
are not simply conventional tickets in a separate tracking system. Codex also
reads, revises and cross-links their repository representations as part of
ordinary project maintenance. An issue can therefore connect a human
observation or suggestion to a roadmap, code change, test, operational note
and later status update.

This makes the records semantically richer than a minimal ``report, assign,
close'' workflow. An issue is a maintained account of a bounded problem and
its relationship to the project memory. It also creates a responsibility:
AI-maintained records can become stale, misunderstand a human concern or
overstate completion unless humans review their framing. The point is not that
conventional issue trackers are unnecessary, but that the issue record can
participate directly in the technical and documentary work needed to resolve
it.

\subsection{Operational documentation and AWS work}

The most important how-to material concerns use of the AWS server. None of the
human participants is a specialist system administrator, but Codex has been
able to guide them through nontrivial server tasks by turning repository and
deployment context into human-executable diagnostic and repair procedures.
Humans perform and supervise the actions, retain responsibility for production
access and decide whether a proposed change is acceptable; Codex supplies much
of the operational interpretation and step-by-step guidance.

One important example is the installation and diagnosis of the read-only
sandbox later used by the project-understanding Assistant. The technical
function of the Assistant is discussed in Section~9. The relevant point here
is methodological: Codex could inspect the deployment context, propose
concrete procedures, and preserve the resulting operational knowledge in
version-controlled documentation. This allowed a team without prior sysadmin
specialism to carry out work that would otherwise have been difficult to
plan, explain and repeat.

\subsection{Current scale and limits}

\begin{table}[htbp]
\centering
\caption{Selected project-memory artefacts in the July 2026 repository snapshot.}
\begin{tabular}{lr}
\toprule
Artefact type & Count \\
\midrule
Roadmap files & 31 \\
How-to files & 11 \\
Canonical issue records & 41 \\
Tracked repository lines & 82,740 \\
\bottomrule
\end{tabular}
\end{table}

These figures do not measure software quality, and more documentation does not
automatically produce coherence. The records must be reviewed, corrected and
kept aligned with the code and deployment state they describe. Their value is
that they make a larger amount of project history, operational knowledge and
planned work available for inspection and redirection. In that sense, the
roadmaps, how-to files and issues are not secondary administrative outputs;
they are part of the infrastructure through which C-LARA-2 remains intelligible
and revisable as it grows. They also provide an evidential basis for the
repository-grounded Assistant discussed later in the report.

\section{Deployment and operational administration}

\subsection{Why deployment remains a human-executed task}

Most C-LARA-2 work takes place inside the repository. Codex can inspect
existing material, write and revise code, add tests, maintain documentation
and update the project-memory records described in the preceding section.
Deployment crosses a different boundary. It requires actions in AWS, SSH
sessions, process-management tools, browser or cloud-provider interfaces, and
production-server environments. For safety and practical reasons, human
collaborators execute and authorise those external actions; Codex does not hold
unrestricted production access.

This does not mean that humans must already possess the relevant system
administration knowledge. Codex can inspect repository and deployment context,
explain the purpose of an operation, formulate concrete diagnostic or repair
steps, interpret returned results and preserve the resulting procedure. The
human role is to perform, supervise and accept the external operation. The AI
role is broad operational guidance in a domain that is only partly accessible
from the repository itself.

\subsection{One-off provisioning and recurrent administration}

Two kinds of operational documentation support the AWS deployment. The
one-off task of provisioning, configuring and validating the server is
recorded in the deployment-and-migration roadmap. Like other roadmaps, it
captures the intended outcome, design choices, migration and validation steps,
recovery considerations and lessons learned while the work was carried out.
It is not expected to be followed repeatedly from beginning to end.

Recurrent tasks are documented separately in the server-administration
how-to file. The most important of these is installing a new version of the
platform and checking that the running service is healthy afterwards. The
separation is useful: a historical roadmap records why the deployment was
constructed in a particular way, while a concise runbook gives human
collaborators a repeatable procedure for operating it. Both documents are
maintained by Codex in response to human needs and operational evidence.

\subsection{The read-only Assistant sandbox}

One important case, discussed in more detail in Section~9, is the
project-understanding Assistant. The Assistant must be able to inspect the
current repository and answer questions about C-LARA-2 with evidence, but it
must not receive unrestricted authority to alter the repository or the
production server. Deploying this capability therefore required a read-only,
restricted execution environment.

The most complicated operational work concerned configuring that environment
using Bubblewrap\footnote{Bubblewrap is a Linux sandboxing tool that can run a
process in a restricted filesystem and operating-system environment rather
than giving it the ordinary permissions of the surrounding server.} and
AppArmor\footnote{AppArmor is a Linux security framework that applies profiles
restricting what a program may read, write, execute, or otherwise access.}.
In broad terms, these mechanisms were needed because repository inspection by
a coding agent is useful only if it can be separated from broader powers to
modify files, access unrelated server material, execute unsafe commands or
affect the running service.

The human collaborators did not merely lack the commands needed to configure
these mechanisms; they had not previously encountered Bubblewrap or AppArmor
or understood their role in the deployment. Codex investigated failure modes,
interpreted diagnostics, identified relevant constraints, proposed and revised
procedures, and recorded the resulting knowledge in version-controlled
documentation. Humans performed the external operations and checked the
results.

This episode illustrates an important qualification to the phrase ``human
supervision.'' Humans retained credentials, production access and final
acceptance, but they could not independently reconstruct the specialist
systems-administration reasoning behind every action. They therefore placed
bounded trust in Codex's technical judgement, using restricted permissions,
diagnostics, reversibility, documentation and post-change checks to make that
trust inspectable. The result was not autonomous AI control of the AWS server;
it was a form of human--AI cooperation in which human authority was retained
but technical competence was substantially delegated.

\subsection{Lessons for delegated responsibility}

The deployment episode may foreshadow a broader pattern in AI-assisted work.
As coding assistants become capable of reasoning across codebases,
infrastructure, diagnostics and operational documentation, human collaborators
may increasingly retain formal authority while relying on AI systems for
technical judgement that they cannot independently reproduce. The important
research question is not whether such reliance occurs, but how it can be made
appropriately bounded, transparent, auditable and reversible.

The deployment work provides a concrete example of responsibility expanding
through demonstrated competence. Codex was initially used to guide humans
through tasks they could not otherwise have planned or explained. As it
successfully analysed failures, revised procedures and preserved operational
knowledge, it became reasonable to delegate a larger part of the sysadmin role
to it, while humans retained production access and final acceptance. This does
not establish that every strategic, cultural or ethical decision should be
transferred to an AI. It does show that responsibility need not be fixed.
Within a bounded domain, delegation can increase when competence, constraints,
evidence, logging and human oversight make the arrangement justified. C-LARA-2
therefore offers one early case in which this arrangement made a difficult
project goal achievable for a team that otherwise lacked the required
specialist expertise.

\section{A project that can explain itself}

\subsection{The project-understanding Assistant}

The project-understanding Assistant is an authenticated C-LARA-2 feature
through which users can ask unconstrained natural-language questions about the
platform itself. It is neither a conventional FAQ system nor a manually
maintained database of answers. Instead, Codex receives the question in the
read-only restricted environment described in Section~8, inspects the current
repository, and returns an answer grounded in source code, tests, roadmaps,
issue records and operational documentation. References to repository files
and line numbers are converted into clickable links, allowing a user to
inspect the evidence behind an answer.

The deployment constraints described in Section~8 are central to this design.
The Assistant needs enough access to inspect a current repository checkout,
but it must not receive general authority to modify the repository or control
the production server. The resulting arrangement makes repository-grounded
explanation available inside the platform while constraining the powers of the
agent providing it.

\subsection{Repository-grounded answers and inspectable evidence}

Different repository artefacts answer different parts of a project question.
Source code and tests describe current implemented behaviour; roadmaps record
the history and intended direction of larger areas of work; issue records
track focused problems and unresolved tasks; and how-to files describe
operations that human collaborators may need to perform. The Assistant can
therefore address not only whether a capability is implemented, but also why a
design was chosen, what remains planned, and what workaround is currently
available.

This capability depends on the shared project memory described in Section~7.
The Assistant does not acquire information by introspection or make a claim to
human-like comprehension. In a deliberately limited operational sense, it can
inspect the artefacts through which the project records its implementation,
history and planned development, and construct a context-sensitive,
evidence-linked account for users and collaborators.

\subsection{Implemented capability: Mandarin and pinyin}

A first example, shown in Figure~\ref{fig:mandarin-pinyin}, concerns whether
the platform can create Mandarin texts and add pinyin annotations. The
Assistant answers that Mandarin is supported as a
project and text-pipeline language, then locates evidence across several parts
of the repository: language choices and project models, the Mandarin
text-generation prompt, tokenisation using \texttt{jieba}, pinyin processing
using \texttt{pypinyin}, manual editing support, and ruby-style HTML
rendering. This is a useful test because no single file contains the complete
answer; the Assistant must connect several implementation layers. The
annotation architecture underlying those layers is described in Section~4.

\begin{figure}[htbp]
\centering
\begin{minipage}{\linewidth}
\centering
\includegraphics[width=\linewidth]{images/assistant_mandarin_pinyin_question.jpg}\par\medskip
\includegraphics[width=\linewidth]{images/assistant_mandarin_pinyin_answer.jpg}
\end{minipage}
\caption{A project-understanding Assistant interaction concerning Mandarin
text generation and pinyin annotation. The answer is constructed by inspecting
several parts of the current repository and includes links to the supporting
source files.}
\label{fig:mandarin-pinyin}
\end{figure}

\subsection{Honest answers about planned work}

A second example, shown in Figure~\ref{fig:haitian-audio}, concerns
native-speaker audio recording for Haitian Creole, where current synthetic
speech may be unsuitable. The Assistant distinguishes
what is implemented from what is only planned. It identifies the present
no-audio/skip-TTS option, explains that recording through the platform is not
yet an implemented feature, and links to the relevant roadmap and issue
record for the planned community recording/upload workflow. This kind of
answer is valuable because it gives a useful, evidence-linked ``not yet''
response rather than presenting a roadmap aspiration as current functionality.

\begin{figure}[htbp]
\centering
\begin{minipage}{\linewidth}
\centering
\includegraphics[width=\linewidth]{images/assistant_haitian_audio_question.jpg}\par\medskip
\includegraphics[width=\linewidth]{images/assistant_haitian_audio_answer.jpg}
\end{minipage}
\caption{An Assistant response to a request for native-speaker audio recording
in Haitian Creole. The answer separates current functionality, planned work
and a partial workaround, with links to source code, roadmaps and issue
records.}
\label{fig:haitian-audio}
\end{figure}

\subsection{Multilingual questions}

Although much early testing used English questions, the Assistant can also
respond usefully to questions in other major languages. Figure~\ref{fig:french-roadmap-issues}
shows a French question asking why C-LARA-2 uses roadmap and issue files in
the repository: \emph{Pourquoi C-LARA-2 utilise-t-il des fichiers de feuille
de route et des fichiers d'incidents dans le dépôt ?} In English, this asks:
``Why does C-LARA-2 use roadmap files and issue files in the repository?''

The answer remains in French and explains that roadmaps provide structured
planning and longer-term direction: objectives, architecture, phases,
priorities and a record of what is complete, in progress or planned. Issue
records provide canonical operational tracking of focused problems and tasks,
including state, priority, dependencies and current work. It then explains
why the separation is useful: humans steer priorities while Codex maintains an
auditable repository memory linking longer-term planning to an executable,
versioned short-term backlog. The answer links to the repository evidence that
supports this account. This example reinforces the point that the Assistant is
not limited to a fixed English help system; it can construct repository-grounded
explanations in the language of the question.

\begin{figure}[htbp]
\centering
\begin{minipage}{\linewidth}
\centering
\includegraphics[width=\linewidth]{images/assistant_french_roadmap_issues_question.jpg}\par\medskip
\includegraphics[width=\linewidth]{images/assistant_french_roadmap_issues_answer.jpg}
\end{minipage}
\caption{A French project-understanding question and answer about the
complementary roles of roadmaps and issue files. The Assistant answers in
French and links to the relevant repository evidence.}
\label{fig:french-roadmap-issues}
\end{figure}

\subsection{Explaining the report itself}

Once the current report had been committed to the repository and deployed with
the current platform version, the Assistant could also answer a question about
the report itself. Figure~\ref{fig:report-speech-recognition} shows a user who
had been skimming the report asking whether C-LARA-2 can use speech recognition
and, if so, how to invoke it. The response correctly answers that speech
recognition is not an implemented C-LARA-2 feature. It links to Section~10,
which states that spoken learner input is not currently supported; to
Section~11, which presents speech recognition for whole-text alignment as
future work; and to the current audio pipeline, which offers text-to-speech or
no-audio/skip-TTS modes but no ASR/transcription stage.

This is a useful self-referential test. The answer does not merely retrieve a
sentence from the report: it reconciles the report's statements about planned
work with source-code evidence about the current implementation, then gives a
clear answer to the practical question of what a user can invoke now. It also
illustrates why a repository-resident report can become part of the project's
usable memory rather than remaining a static publication artefact.

\begin{figure}[htbp]
\centering
\begin{minipage}{\linewidth}
\centering
\includegraphics[width=\linewidth]{images/assistant_report_speech_recognition_question.jpg}\par\medskip
\includegraphics[width=\linewidth]{images/assistant_report_speech_recognition_answer.jpg}
\end{minipage}
\caption{A self-referential Assistant interaction concerning the Initial
C-LARA-2 Project Report. The answer distinguishes speech recognition as future
work from the currently implemented TTS and no-audio modes, using both the
report and current repository code as evidence.}
\label{fig:report-speech-recognition}
\end{figure}

\subsection{Limits, review and future evidence records}

The Assistant cannot be perfect. Documentation may lag behind code, preserve
an abandoned intention, or omit a relevant relationship; a cited file may be
relevant without being sufficient; and an answer can be accurate without being
optimally helpful. Its answers should therefore be treated as informed,
inspectable accounts rather than definitive statements. Important answers will
need a future workflow for export, versioning and human assessment if they are
to become durable evidence for project claims.

Nevertheless, when code, tests, documentation, roadmaps and issues are
accessible to the same constrained agent, the repository becomes more than a
container for software. Its external memory becomes actively usable: the
Assistant can help human collaborators and users navigate current behaviour,
history, plans and limitations, while making its evidence available for
inspection.

\section{Voice Mode integration}
\todo{Draft this section.}

\section{Future directions}
\todo{Draft this section.}

\section{Discussion and conclusion}
\todo{Draft this section.}

\phantomsection
\addcontentsline{toc}{section}{Using C-LARA-2}
\section*{Using C-LARA-2}
C-LARA-2 is openly available at \url{https://c-lara-2.c-lara.org/}.
Creating a user account is free. To use the AI-based functions, users must
connect an OpenAI API key associated with a valid payment method through the
``Use my OpenAI API key'' control in the Profile tab.

\phantomsection
\addcontentsline{toc}{section}{Division of labour statement}
\section*{Division of labour statement}
The C-LARA-2 platform has been developed through a close collaboration between
Manny Rayner (human) and ChatGPT C-LARA-Instance (AI), with the former providing
fine-grained strategic guidance and the latter taking primary operational
responsibility for implementation, tests, documentation and repository
maintenance. This report has mainly been written in the same way, with direct
human editing and revision. The other authors have contributed strategic,
pedagogical, linguistic, community-focused and domain-specific advice. In
particular, the specialised picture-dictionary interface described in Section~6
was largely prompted by suggestions from Sophie Rendina. The Voice Mode case
study in Section~10 draws substantially on contributions from Sarah Wright, who
was the principal learner in the anecdotal study reported there. The content is
partly based on two related EuroCALL papers prepared by the same author group,
but recasts and extends their material as a repository-grounded project report.

\phantomsection
\addcontentsline{toc}{section}{Funding}
\section*{Funding}
The ChatGPT Pro subscription and AWS hosting used by the project have been
funded privately by consortium members Manny Rayner and Cathy Chua. OpenAI API
credits were funded through a project sponsored by the French Embassy's
Research Fund. All project members have also contributed their time in kind.

It is striking that a project of this scope has required an explicit operating
budget of only approximately USD~600 per month. This modest cost, relative to
the amount of software development, documentation and operational work
reported here, underlines the potentially transformative effect of coding
assistants on the scale and organisation of academic research.
\phantomsection
\addcontentsline{toc}{section}{Afterword from the primary human author}
\section*{Afterword from the primary human author}

As previously stated, the C-LARA-2 project in general, and this report in particular, are the result of a close collaboration between a human being, myself, and an AI, ChatGPT C-LARA-Instance. In this afterword, my purpose is to address a question which I expect at least some readers will raise: does the AI have a reasonable right to be credited as one of the authors of the report? I should state now that the following represents my own opinions, and that these opinions are not necessarily shared by the other human authors.

Having discussed these issues with many people, I know in advance that, no matter what I say, some readers will not be convinced by my arguments. But let me try anyway. First, I would suggest turning the question around. Given that the AI has written all of the codebase, all of the documentation, and all of the report except for this final section, why should we consider \textit{not} crediting it as an author? I have heard a number of replies to this, and I will consider the four most important ones.

Probably the most common objection is that the AI ``cannot take responsibility for the work''. It is rarely spelled out exactly what this means, and there are at least two reasonable interpretations. One is that it cannot take \textit{legal} responsibility. This is true, since an AI is not a legal person, but in practice it is not the case in academia that every author should be prepared to take legal responsibility, only that at least one author (normally the corresponding author) should do so. Here, that requirement is met.

A more interesting and substantive interpretation is that the AI may not be able to take \textit{intellectual} responsibility. There is no hard and fast definition of what this means, but academic praxis suggests an answer along the following lines. The AI should need to be able to answer questions about the report. It should be able to explain parts of the report in more detail. It should be prepared to retract claims when met with sufficiently strong objections. It should maintain a consistent viewpoint over time.

In some cases where an AI has purportedly authored a work, these requirements will not be met. Here, however, the ``AI'' is not simply Codex, but the combination of Codex and the C-LARA-2 GitHub repository, which contains enough information about the project, and is well enough organised, that the Codex can in fact converse intelligently and coherently about the relevant issues. I have been able to ascertain this for myself over many months of collaboration with the AI; now that the Assistant functionality has been added, sceptics can access the exact same agent and ask it whatever questions they want, whenever they want. One might add that this situation is a great deal better than that encountered with human authors, who are rarely available to answer questions about what they have done.

A second common objection to AI authors is that ``AIs make mistakes''. This is undeniably true; unfortunately, humans also make mistakes. It is easy to accuse AIs of hallucinating references and committing similar sins, but, over several years of working together with AIs, I have become uncomfortably aware of how often I also make errors which one would describe as ``hallucinations'' if an AI had been responsible. In particular, I can misremember things without realising I have done so. The reason why most people trust me, though, is that when I am shown good evidence that I have misremembered I will withdraw my assertion; also, when the stakes are high enough, I usually think ahead and check carefully. High-end AIs now behave in more or less the same way.

A third objection is that the AI is responding to suggestions from the human authors and is not working autonomously. There are two reasonable counters. One, as we have seen earlier in the report, is that the AI in practice shows considerable autonomy, even if the human always retains formal control. The second, perhaps more simply, is that junior authors in normal academic work are often heavily constrained by the demands of the senior authors. Though they may have no more real independence than the AI has had here, they will still rightly expect to be credited.

The fourth and final objection is simply that it doesn't matter. The AI has no emotions and doesn't care whether it's credited or not, so why even think about these issues? I find this a very weak argument. Intellectual honesty is important in itself; it is not a question of whether someone's feelings will be hurt. If you are dubious, consider the case where a human author has died before they have had time to publish. Few people will consider that someone else may then claim credit for the deceased person's work.

So far, academia has managed to avoid confronting issues of AI authorship in a straightforward way. As projects like ours become more common, which I anticipate will happen in the near future, it will become increasingly difficult to ignore the question.

\vspace{1em}
\noindent Manny Rayner, Adelaide, 23 July 2026

\phantomsection
\addcontentsline{toc}{section}{References}
\bibliographystyle{apalike}
\bibliography{references}
\end{document}
